% Obsahový zdroj: import_hodiny/30_Pravidlo_souctu_a_soucinu.tex a odpovídající PDF (obě strany).
% P1–P5: zdrojové Příklady 1–5; C1–C5: zdrojové Procvičování 1–5.
% Matematický obsah zachován; výklad, odvození a řešení pedagogicky rozvedeny.
% Upřesnění a opravy: viz KMB_29-41_Kontrola.md; původní import zůstává beze změny.
% Rozpis: 4_rocnik/tematicke_plany/Rozpis_vyuky_MAT_4_rocnik.tex, hodina 30.
% Plán: 4_rocnik/tematicke_plany/SPS.pdf, oddíly 2.1–2.3.
% Sazba: sablony/pripravy/TEMPLATE_Priprava.tex a společný priprava.sty.
% Vlastní doplnění: didaktický výklad, zdůvodnění, metodika; žádné nové číslované úlohy.
% Učebnice/09_Kombinatorika nemá obsahovou kapitolu; místní kombinatorické listy ani Realisticky.cz nejsou dostupné.
\documentclass[11pt,a4paper]{article}
% Spolecny styl pro ucitelske pripravy MAT (Work/MAT)
% Verze 1.1, 2026-10-23
\NeedsTeXFormat{LaTeX2e}
\ProvidesPackage{priprava}[2026/10/23 v1.1 Sablona ucitelskych priprav MAT]

% Zakladni typografie a cestina
\RequirePackage[T1]{fontenc}
\RequirePackage[utf8]{inputenc}
\RequirePackage[czech]{babel}
\RequirePackage{lmodern}
\RequirePackage{microtype}
\RequirePackage{amsmath,amssymb}
\RequirePackage{xcolor}
\RequirePackage{enumitem}
\RequirePackage[most]{tcolorbox}
\RequirePackage{tabularx,booktabs,array}
\RequirePackage{fancyhdr}
\RequirePackage{needspace}
\RequirePackage{geometry}
\RequirePackage[hidelinks]{hyperref}

\geometry{a4paper,top=17mm,bottom=18mm,left=23mm,right=23mm}
\definecolor{notegray}{gray}{0.39}
\setlength{\parindent}{0pt}
\setlength{\parskip}{2.5pt plus 1pt minus 0.5pt}
\setlist[itemize]{topsep=2pt,itemsep=1pt,parsep=0pt,leftmargin=1.45em}
\setlist[enumerate]{topsep=2pt,itemsep=2pt,parsep=0pt,leftmargin=1.65em}
\setcounter{secnumdepth}{0}
\renewcommand{\arraystretch}{1.12}

% Zahlavi: vlevo rocnik | globalni tema, vpravo cislo: nazev hodiny
% Zapati zustava zcela prazdne.
\pagestyle{fancy}
\fancyhf{}
\renewcommand{\headrulewidth}{0.35pt}
\renewcommand{\footrulewidth}{0pt}
\setlength{\headheight}{14pt}
\newcommand{\nastavPripravu}[4]{%
  \fancyhead[L]{\small #1\enspace\textbar\enspace #2}%
  \fancyhead[R]{\small #3: #4}%
}

% Minimalisticke cernobile boxy. Kazda definice zvlast.
\tcbset{commonlesson/.style={enhanced,breakable=false,colback=white,colframe=black,
  colbacktitle=white,coltitle=black,boxrule=0.45pt,
  arc=0.8mm,outer arc=0.8mm,left=3mm,right=3mm,top=1.6mm,bottom=1.6mm,
  toptitle=1mm,bottomtitle=0.7mm,fonttitle=\bfseries,
  attach boxed title to top left={xshift=2.5mm,yshift=-1.15mm},
  boxed title style={colback=white,colframe=white,boxrule=0pt,left=1mm,right=1mm}}}
\newtcolorbox{zadaniBox}[1][]{commonlesson,title={Zadání},#1}
\newtcolorbox{prikladBox}[1][]{commonlesson,title={Příklad},#1}
\newtcolorbox{definiceBox}[1][]{commonlesson,title={Definice},#1}

% Spolecny box pro zadani a jeho reseni, bez mezery mezi castmi.
% Pouziti:
%   \begin{prikladReseniBox}
%     Zadani ulohy.
%     \tcblower
%     \begin{align*} ... \end{align*}
%   \end{prikladReseniBox}
% Spodni cast ma jemny sedy podklad a delici linku; lze v ni sazet
% vice radku matematiky. Obe casti maji jen jeden spolecny ramecek.
\tcbset{solutionlesson/.style={
  commonlesson,
  enhanced jigsaw,
  breakable,
  colbacklower=black!3,
  segmentation style={solid,draw=black!25,line width=0.35pt},
  fontlower=\small\color{notegray},
  before lower={\noindent\textbf{Řešení:}\enspace},
  leftlower=3mm,rightlower=3mm
}}
\newtcolorbox{prikladReseniBox}[1][]{solutionlesson,title={Příklad},#1}
\newtcolorbox{zadaniReseniBox}[1][]{solutionlesson,title={Zadání},#1}

% Pomucky pro autora pripravy: sedive poznamky; \reseni pro starsi pripravy
\newcommand{\poznamka}[1]{\par\begingroup\small\color{notegray}\itshape
  \textnormal{\textbf{Poznámka pro učitele.}}\enspace #1\par\endgroup}
\newcommand{\reseni}[1]{\par\begingroup\small\color{notegray}
  \textbf{Očekávané řešení:}\enspace #1\par\endgroup}
\newcommand{\blok}[1]{\needspace{5\baselineskip}\subsection*{#1}}
\newcommand{\vyrok}[1]{\uv{#1}}

\nastavPripravu{4. ročník}{Kombinatorika}{30}{Pravidlo součtu a součinu}
\begin{document}


Vrať se ke dvěma způsobům výpisu z hodiny 29. Ve skupinách alternativ sčítáme počty větví, při postupném vytváření jednoho výsledku zjišťujeme počet všech pokračování.
\poznamka{Předpoklady: pořadí, opakování a úplný výpis. Cíl: žák zdůvodní volbu součtu či součinu a správně započítá omezení na pozicích. Společně řeš P1–P4; P5 nabídni jako rozšíření rychlejším žákům.}

\blok{Výklad a odvození}
\begin{definiceBox}
\textbf{Kombinatorické pravidlo součtu.} Rozdělíme-li všechny přípustné výsledky na vzájemně disjunktní (nepřekrývající se) skupiny o počtech $a_1,\ldots,a_r$, je všech výsledků $a_1+\cdots+a_r$. Každý výsledek musí patřit právě do jedné skupiny.
\end{definiceBox}
\begin{definiceBox}
\textbf{Kombinatorické pravidlo součinu.} Vytváříme-li výsledek v $r$ krocích a po každé přípustné volbě v předchozích krocích má krok $i$ vždy právě $a_i$ pokračování, je výsledků $a_1a_2\cdots a_r$. Každý sled voleb musí jednoznačně určit právě jeden výsledek.
\end{definiceBox}
\begin{definiceBox}
\textbf{Strom možností} je diagram postupných voleb: větve představují povolené volby, úplná cesta od začátku ke konci jeden výsledek.
\end{definiceBox}

U dvou kroků si představ $a$ skupin, každou s $b$ pokračováními. Opakovaný součet $b+\cdots+b$ ($a$ sčítanců) dává $ab$. Stejně zdůvodníme více kroků.

\textbf{Strom pro dvě první cifry trojciferného čísla z P4:}
\[
\begin{array}{c|c|c}
\text{první cifra}&\text{druhá cifra}&\text{poslední cifra}\\\hline
1&2,3,4&0\\
2&1,3,4&0\\
3&1,2,4&0\\
4&1,2,3&0
\end{array}
\]
Každý ze čtyř začátků má tři pokračování a jedno pevné zakončení. Tabulka je úsporným zachycením větví stromu, nikoli seznamem hotových čísel.

Dovolené znaky mohou záviset na minulých volbách; pro součin je podstatný \emph{stejný počet} pokračování v daném kroku. Není nutné, aby zůstávala stejná množina znaků. Liší-li se i počty, rozdělujeme na případy se stejným větvením a výsledky sčítáme.


\poznamka{Průběh: 0–6 min návrat k výpisu; 6–17 min pravidla a P1–P2; 17–30 min P3–P4, strom; 30–41 min C1–C4, při dostatku času C5; 41–45 min kontrola. P5 a případně C5 jsou rezerva či domácí práce, aby neubraly čas vysvětlení počáteční nuly.}

\blok{Společně řešené příklady}

% Zdroj P1: import_hodiny/30_Pravidlo_souctu_a_soucinu.tex, Příklad 1.
\begin{prikladReseniBox}[title={Příklad P1}]
Jízdenku lze koupit ve čtyřech různých tarifech na vlak nebo ve třech tarifech na autobus. Kolik je možností jedné jízdenky?
\tcblower
Volíme buď vlak, nebo autobus, nikoli oba zároveň. Množiny typů jízdenek jsou disjunktní. \(4+3=7\).
\end{prikladReseniBox}


% Zdroj P2: import_hodiny/30_Pravidlo_souctu_a_soucinu.tex, Příklad 2.
\begin{prikladReseniBox}[title={Příklad P2}]
Na oběd je pět hlavních jídel a tři nápoje. Vybíráme jedno z každého.
\tcblower
Nejprve \(5\) možností jídla a pro každé z nich \(3\) možnosti nápoje. \(5\cdot3=15\).
\end{prikladReseniBox}


% Zdroj P3: import_hodiny/30_Pravidlo_souctu_a_soucinu.tex, Příklad 3.
\begin{prikladReseniBox}[title={Příklad P3}]
Kolik existuje čtyřciferných čísel bez opakování číslic?
\tcblower
Obsazujeme postupně čtyři pozice. První cifra je z $1,\ldots,9$: devět voleb. Pak zbývá devět nepoužitých číslic, protože se vrací možnost nuly. Další pozice mají osm a sedm možností. Pravidlem součinu
\[
9\cdot9\cdot8\cdot7=4536.
\]
Každé čtyřciferné číslo má právě jeden takový zápis. Součin $10\cdot9\cdot8\cdot7$ by navíc započítal kódy začínající nulou.
\end{prikladReseniBox}


% Zdroj P4: import_hodiny/30_Pravidlo_souctu_a_soucinu.tex, Příklad 4.
\begin{prikladReseniBox}[title={Příklad P4}]
Kolik trojciferných čísel s navzájem různými číslicemi dělitelných pěti vytvoříme z číslic $0,1,2,3,4$?
\tcblower
Dělitelnost pěti vyžaduje poslední cifru $0$ nebo $5$. Zadaná nabídka neobsahuje $5$, takže konec je pevně $0$. Na prvním místě jsou čtyři nenulové možnosti, na druhém tři zbývající nenulové. Dostáváme $4\cdot3=12$ čísel. Například větev začínající $1$ dává $120,130,140$; obdobně další tři větve.
\end{prikladReseniBox}


% Zdroj P5: import_hodiny/30_Pravidlo_souctu_a_soucinu.tex, Příklad 5.
\begin{prikladReseniBox}[title={Příklad P5}]
Kolik čtyřznakových kódů z $A,B,C,D$ lze sestavit, nesmějí-li být stejné dva sousední znaky? Opakování na nesousedních pozicích je dovoleno.
\tcblower
První místo má čtyři volby. Na druhém jsou tři (vše kromě prvního znaku). Na třetím opět tři, protože vylučujeme jen druhý znak, a na čtvrtém zase tři. Celkem $4\cdot3\cdot3\cdot3=108$. Množina dovolených znaků se mění, její počet ale zůstává stejný, a proto lze použít součin.
\end{prikladReseniBox}


\newpage
\blok{Samostatná práce a zásoba k procvičení}

U každé úlohy nejprve žák napíše, zda volí alternativu, nebo sestavuje výsledek v několika krocích. C1–C2 tvoří krátký kontrast; těžiště samostatné práce je v C3–C4.

% Zdroj C1: import_hodiny/30_Pravidlo_souctu_a_soucinu.tex, Procvičování, položka 1.
\begin{zadaniReseniBox}[title={Úloha C1}]
Tři různé příchutě a čtyři polevy: jedna příchuť a jedna poleva.
\tcblower
Pro každou ze tří příchutí jsou čtyři polevy. Jeden výsledek je dvojice příchuť–poleva, tedy $3\cdot4=12$ možností.
\end{zadaniReseniBox}

% Zdroj C2: import_hodiny/30_Pravidlo_souctu_a_soucinu.tex, Procvičování, položka 2.
\begin{zadaniReseniBox}[title={Úloha C2}]
Do finále se lze dostat jednou z pěti postupových cest A nebo jednou ze sedmi cest B; cesty se nepřekrývají.
\tcblower
Vybereme jednu cestu z jedné ze dvou disjunktních skupin. Sčítáme $5+7=12$ cest, protože nejde o dvě postupné volby.
\end{zadaniReseniBox}

% Zdroj C3: import_hodiny/30_Pravidlo_souctu_a_soucinu.tex, Procvičování, položka 3.
\begin{zadaniReseniBox}[title={Úloha C3}]
Kolik čtyřmístných kódů z číslic \(0\) až \(5\) lze napsat, je-li opakování dovoleno?
\tcblower
Na každé ze čtyř pozic je šest možností $0,\ldots,5$. Opakování ani počáteční nula nejsou zakázány, tedy $6\cdot6\cdot6\cdot6=1296$ kódů.
\end{zadaniReseniBox}

% Zdroj C4: import_hodiny/30_Pravidlo_souctu_a_soucinu.tex, Procvičování, položka 4.
\begin{zadaniReseniBox}[title={Úloha C4}]
Kolik trojciferných čísel lze napsat z číslic \(0\) až \(6\), bez opakování?
\tcblower
První cifra: šest možností $1,\ldots,6$. Druhá: šest zbývajících včetně nuly. Třetí: pět zbývajících. Počet je $6\cdot6\cdot5=180$.
\end{zadaniReseniBox}

% Zdroj C5: import_hodiny/30_Pravidlo_souctu_a_soucinu.tex, Procvičování, položka 5.
\begin{zadaniReseniBox}[title={Úloha C5}]
Kolik existuje pětimístných hesel ze tří znaků, v nichž žádné dva sousední nejsou stejné?
\tcblower
První znak má tři možnosti. Na každé z dalších čtyř pozic vyloučíme pouze bezprostředně předcházející znak, a proto zbývají vždy dvě volby. Celkem $3\cdot2^4=48$ hesel.
\end{zadaniReseniBox}


\poznamka{Typické chyby: sčítání počtů jídel a nápojů; překrývající se případy; druhé místo v P3 mylně jen s osmi možnostmi. Po první nenulové cifře zbývá devět cifer včetně nuly. U P5 je zakázáno pouze sousední opakování, nikoli návrat znaku později (např. $ABAB$). Příště zavedeme faktoriál jako zápis zkracující dlouhé součiny.}
\end{document}
